% --- Archon physics formalization source begin ---
% archon:physics
% archon:covers PhyXMiniProblems/problem_phyx_mini_0407.lean
% archon:source-report reports/phyx_mini/problem_phyx_mini_0407.source.json

\chapter{Physics problem phyx\_mini\_0407}
\label{ch:PhyXMiniProblems_problem_phyx_mini_0407}

\paragraph{Problem source.}
\#\# Physical scenario

A 6.0-\$	ext\{cm\}\$-diameter cylinder of nitrogen gas has a 4.0-\$	ext\{cm\}\$-thick movable copper piston. The cylinder is oriented vertically, as shown in figure, and the air above the piston is evacuated. When the gas temperature is 20 \$\textasciicircum{}\{\textbackslash{}circ\}\textbackslash{}text\{C\}\$, the piston floats 20 \$	ext\{cm\}\$ above the bottom of the cylinder.

\#\# Question

How much work is done on the gas as the piston rises?

\#\# Answer choices

- A: A: 0.11 J

- B: B: 1.2 J

- C: C: 0.57 J

- D: D: -0.57 J

\#\# Figure caption (auxiliary; use the image as primary evidence)

The image depicts a cross-sectional view of a vertical cylinder containing a piston. Here are the elements present:

- **Cylinder**: A vertical container, outlined with borders, possibly indicating the wall of the cylinder.

- **Piston**: A horizontal gray bar inside the cylinder, separating two regions. It is labeled "Piston."

- **Vacuum Area**: The top section above the piston is labeled "Vacuum," indicating no matter present in this space.

- **Gas (N₂)**: Below the piston, a yellow section is labeled "N₂," representing nitrogen gas.

- **Dimensions**: 
  - The vertical distance from the top to the bottom of the cylinder is indicated as "20 cm" with a double-headed arrow pointing downward.
  - The base of the section containing N₂ is labeled with a width of "6.0 cm."

This setup likely represents a physics or thermodynamics concept involving pressure, volume, or gas laws.

\#\# Dataset metadata

Ideal Gases and Kinetic Theory; Physical Model Grounding Reasoning, Multi-Formula Reasoning

\paragraph{Recorded answer/context.}
C: C: 0.57 J

\paragraph{Figure/image path.}
/root/proposal\_for\_physic/science-mango/phyx\_mini\_run/phyx\_data/test\_image/407.png

\paragraph{Formalization target.}
create a compiling Lean file with sorry bodies at `PhyXMiniProblems/problem\_phyx\_mini\_0407.lean`. The Lean declarations must preserve the physical quantities, dimensions or dimensional roles, figure labels, governing-law hypotheses, and final relation expressed by this problem.
Use Mathlib/Physlib names found through LeanExplore where available. If a physics API is missing, introduce faithful local abstractions rather than scalar placeholder aliases.

\begin{theorem}[Physics formalization target]
\label{thm:physics:phyx_mini_0407:target}
\lean{PhyXMiniProblems.ProblemPhyXMini0407.recorded_positive_choice_C_conflicts_with_rising_work_on_gas}
\uses{def:physics:phyx-mini-0407:phyxminiproblems-problemphyxmini0407-lengthinmeters, def:physics:phyx-mini-0407:phyxminiproblems-problemphyxmini0407-copperpistonnitrogencylinder, def:physics:phyx-mini-0407:phyxminiproblems-problemphyxmini0407-matchesproblemandprimaryfigure, def:physics:phyx-mini-0407:phyxminiproblems-problemphyxmini0407-hasphysicalparameters, def:physics:phyx-mini-0407:phyxminiproblems-problemphyxmini0407-obeyspistoncylinderlaws, def:physics:phyx-mini-0407:phyxminiproblems-problemphyxmini0407-workongasmatcheschoice}
The assigned autoformalize agent should translate this physics problem into Lean declarations in the covered file, with theorem and lemma proof bodies written as `by sorry`.
\end{theorem}
\begin{proof}
This is an autoformalization task, not a proof task. Produce statements that can later be proved by the physics prover without weakening the physical model.
\end{proof}

% --- Archon declaration topology begin ---
\section{Declaration topology}
\begin{definition}[Length Quantity]
  \label{def:physics:phyx-mini-0407:phyxminiproblems-problemphyxmini0407-lengthquantity}
  \lean{PhyXMiniProblems.ProblemPhyXMini0407.LengthQuantity}
  A physical length, independent of the unit used to read it
\end{definition}
\begin{proof}
  This is the declared model definition of Length Quantity.
\end{proof}
\begin{definition}[Mass Quantity]
  \label{def:physics:phyx-mini-0407:phyxminiproblems-problemphyxmini0407-massquantity}
  \lean{PhyXMiniProblems.ProblemPhyXMini0407.MassQuantity}
  A physical mass, independent of the unit used to read it
\end{definition}
\begin{proof}
  This is the declared model definition of Mass Quantity.
\end{proof}
\begin{definition}[Mass Density Quantity]
  \label{def:physics:phyx-mini-0407:phyxminiproblems-problemphyxmini0407-massdensityquantity}
  \lean{PhyXMiniProblems.ProblemPhyXMini0407.MassDensityQuantity}
  A physical mass density, with dimension mass per volume
\end{definition}
\begin{proof}
  This is the declared model definition of Mass Density Quantity.
\end{proof}
\begin{definition}[Volume Quantity]
  \label{def:physics:phyx-mini-0407:phyxminiproblems-problemphyxmini0407-volumequantity}
  \lean{PhyXMiniProblems.ProblemPhyXMini0407.VolumeQuantity}
  A physical volume, with length-cubed dimension
\end{definition}
\begin{proof}
  This is the declared model definition of Volume Quantity.
\end{proof}
\begin{definition}[Acceleration Quantity]
  \label{def:physics:phyx-mini-0407:phyxminiproblems-problemphyxmini0407-accelerationquantity}
  \lean{PhyXMiniProblems.ProblemPhyXMini0407.AccelerationQuantity}
  A physical acceleration, used for local gravitational acceleration
\end{definition}
\begin{proof}
  This is the declared model definition of Acceleration Quantity.
\end{proof}
\begin{definition}[Temperature Quantity]
  \label{def:physics:phyx-mini-0407:phyxminiproblems-problemphyxmini0407-temperaturequantity}
  \lean{PhyXMiniProblems.ProblemPhyXMini0407.TemperatureQuantity}
  An absolute thermodynamic temperature
\end{definition}
\begin{proof}
  This is the declared model definition of Temperature Quantity.
\end{proof}
\begin{definition}[Molar Gas Constant Quantity]
  \label{def:physics:phyx-mini-0407:phyxminiproblems-problemphyxmini0407-molargasconstantquantity}
  \lean{PhyXMiniProblems.ProblemPhyXMini0407.MolarGasConstantQuantity}
  A molar gas constant, omitting only Physlib's unavailable mole dimension
\end{definition}
\begin{proof}
  This is the declared model definition of Molar Gas Constant Quantity.
\end{proof}
\begin{definition}[length In Meters]
  \label{def:physics:phyx-mini-0407:phyxminiproblems-problemphyxmini0407-lengthinmeters}
  \lean{PhyXMiniProblems.ProblemPhyXMini0407.lengthInMeters}
  \uses{def:physics:phyx-mini-0407:phyxminiproblems-problemphyxmini0407-lengthquantity}
  Metre readout of a physical length
\end{definition}
\begin{proof}
  This is the declared model definition of length In Meters.
\end{proof}
\begin{definition}[mass In Kilograms]
  \label{def:physics:phyx-mini-0407:phyxminiproblems-problemphyxmini0407-massinkilograms}
  \lean{PhyXMiniProblems.ProblemPhyXMini0407.massInKilograms}
  \uses{def:physics:phyx-mini-0407:phyxminiproblems-problemphyxmini0407-massquantity}
  Kilogram readout of a physical mass
\end{definition}
\begin{proof}
  This is the declared model definition of mass In Kilograms.
\end{proof}
\begin{definition}[mass Density In Kilograms Per Cubic Meter]
  \label{def:physics:phyx-mini-0407:phyxminiproblems-problemphyxmini0407-massdensityinkilogramspercubicmeter}
  \lean{PhyXMiniProblems.ProblemPhyXMini0407.massDensityInKilogramsPerCubicMeter}
  \uses{def:physics:phyx-mini-0407:phyxminiproblems-problemphyxmini0407-massdensityquantity}
  Kilogram-per-cubic-metre readout of a physical mass density
\end{definition}
\begin{proof}
  This is the declared model definition of mass Density In Kilograms Per Cubic Meter.
\end{proof}
\begin{definition}[area In Square Meters]
  \label{def:physics:phyx-mini-0407:phyxminiproblems-problemphyxmini0407-areainsquaremeters}
  \lean{PhyXMiniProblems.ProblemPhyXMini0407.areaInSquareMeters}
  Square-metre readout of a nonnegative physical area
\end{definition}
\begin{proof}
  This is the declared model definition of area In Square Meters.
\end{proof}
\begin{definition}[volume In Cubic Meters]
  \label{def:physics:phyx-mini-0407:phyxminiproblems-problemphyxmini0407-volumeincubicmeters}
  \lean{PhyXMiniProblems.ProblemPhyXMini0407.volumeInCubicMeters}
  \uses{def:physics:phyx-mini-0407:phyxminiproblems-problemphyxmini0407-volumequantity}
  Cubic-metre readout of a physical volume
\end{definition}
\begin{proof}
  This is the declared model definition of volume In Cubic Meters.
\end{proof}
\begin{definition}[pressure In Pascals]
  \label{def:physics:phyx-mini-0407:phyxminiproblems-problemphyxmini0407-pressureinpascals}
  \lean{PhyXMiniProblems.ProblemPhyXMini0407.pressureInPascals}
  Pascal readout of a physical pressure
\end{definition}
\begin{proof}
  This is the declared model definition of pressure In Pascals.
\end{proof}
\begin{definition}[acceleration In Meters Per Second Squared]
  \label{def:physics:phyx-mini-0407:phyxminiproblems-problemphyxmini0407-accelerationinmeterspersecondsquared}
  \lean{PhyXMiniProblems.ProblemPhyXMini0407.accelerationInMetersPerSecondSquared}
  \uses{def:physics:phyx-mini-0407:phyxminiproblems-problemphyxmini0407-accelerationquantity}
  Metres-per-second-squared readout of a physical acceleration
\end{definition}
\begin{proof}
  This is the declared model definition of acceleration In Meters Per Second Squared.
\end{proof}
\begin{definition}[temperature In Kelvins]
  \label{def:physics:phyx-mini-0407:phyxminiproblems-problemphyxmini0407-temperatureinkelvins}
  \lean{PhyXMiniProblems.ProblemPhyXMini0407.temperatureInKelvins}
  \uses{def:physics:phyx-mini-0407:phyxminiproblems-problemphyxmini0407-temperaturequantity}
  Kelvin readout of an absolute thermodynamic temperature
\end{definition}
\begin{proof}
  This is the declared model definition of temperature In Kelvins.
\end{proof}
\begin{definition}[temperature In Degrees Celsius]
  \label{def:physics:phyx-mini-0407:phyxminiproblems-problemphyxmini0407-temperatureindegreescelsius}
  \lean{PhyXMiniProblems.ProblemPhyXMini0407.temperatureInDegreesCelsius}
  \uses{def:physics:phyx-mini-0407:phyxminiproblems-problemphyxmini0407-temperaturequantity, def:physics:phyx-mini-0407:phyxminiproblems-problemphyxmini0407-temperatureinkelvins}
  Celsius readout obtained from the kelvin readout by the standard offset
\end{definition}
\begin{proof}
  This is the declared model definition of temperature In Degrees Celsius.
\end{proof}
\begin{definition}[molar Gas Constant In SI]
  \label{def:physics:phyx-mini-0407:phyxminiproblems-problemphyxmini0407-molargasconstantinsi}
  \lean{PhyXMiniProblems.ProblemPhyXMini0407.molarGasConstantInSI}
  \uses{def:physics:phyx-mini-0407:phyxminiproblems-problemphyxmini0407-molargasconstantquantity}
  SI readout of the molar gas constant in joules per mole-kelvin
\end{definition}
\begin{proof}
  This is the declared model definition of molar Gas Constant In SI.
\end{proof}
\begin{definition}[energy In Joules]
  \label{def:physics:phyx-mini-0407:phyxminiproblems-problemphyxmini0407-energyinjoules}
  \lean{PhyXMiniProblems.ProblemPhyXMini0407.energyInJoules}
  Joule readout of a signed physical energy
\end{definition}
\begin{proof}
  This is the declared model definition of energy In Joules.
\end{proof}
\begin{definition}[Gas Species]
  \label{def:physics:phyx-mini-0407:phyxminiproblems-problemphyxmini0407-gasspecies}
  \lean{PhyXMiniProblems.ProblemPhyXMini0407.GasSpecies}
  The gas species stated in the problem and printed in the figure
\end{definition}
\begin{proof}
  This is the declared model definition of Gas Species.
\end{proof}
\begin{definition}[Piston Material]
  \label{def:physics:phyx-mini-0407:phyxminiproblems-problemphyxmini0407-pistonmaterial}
  \lean{PhyXMiniProblems.ProblemPhyXMini0407.PistonMaterial}
  The material of the movable piston
\end{definition}
\begin{proof}
  This is the declared model definition of Piston Material.
\end{proof}
\begin{definition}[Cylinder Orientation]
  \label{def:physics:phyx-mini-0407:phyxminiproblems-problemphyxmini0407-cylinderorientation}
  \lean{PhyXMiniProblems.ProblemPhyXMini0407.CylinderOrientation}
  Orientation of the cylindrical apparatus
\end{definition}
\begin{proof}
  This is the declared model definition of Cylinder Orientation.
\end{proof}
\begin{definition}[Above Piston Medium]
  \label{def:physics:phyx-mini-0407:phyxminiproblems-problemphyxmini0407-abovepistonmedium}
  \lean{PhyXMiniProblems.ProblemPhyXMini0407.AbovePistonMedium}
  Medium occupying the region above the piston
\end{definition}
\begin{proof}
  This is the declared model definition of Above Piston Medium.
\end{proof}
\begin{definition}[Piston Mobility]
  \label{def:physics:phyx-mini-0407:phyxminiproblems-problemphyxmini0407-pistonmobility}
  \lean{PhyXMiniProblems.ProblemPhyXMini0407.PistonMobility}
  Kinematic role of the piston stated in the prose
\end{definition}
\begin{proof}
  This is the declared model definition of Piston Mobility.
\end{proof}
\begin{definition}[Vertical Motion]
  \label{def:physics:phyx-mini-0407:phyxminiproblems-problemphyxmini0407-verticalmotion}
  \lean{PhyXMiniProblems.ProblemPhyXMini0407.VerticalMotion}
  Direction of the process named in the question
\end{definition}
\begin{proof}
  This is the declared model definition of Vertical Motion.
\end{proof}
\begin{definition}[Figure Region Label]
  \label{def:physics:phyx-mini-0407:phyxminiproblems-problemphyxmini0407-figureregionlabel}
  \lean{PhyXMiniProblems.ProblemPhyXMini0407.FigureRegionLabel}
  Text labels visible in the primary raster
\end{definition}
\begin{proof}
  This is the declared model definition of Figure Region Label.
\end{proof}
\begin{definition}[Copper Piston Nitrogen Cylinder]
  \label{def:physics:phyx-mini-0407:phyxminiproblems-problemphyxmini0407-copperpistonnitrogencylinder}
  \lean{PhyXMiniProblems.ProblemPhyXMini0407.CopperPistonNitrogenCylinder}
  \uses{def:physics:phyx-mini-0407:phyxminiproblems-problemphyxmini0407-lengthquantity, def:physics:phyx-mini-0407:phyxminiproblems-problemphyxmini0407-massquantity, def:physics:phyx-mini-0407:phyxminiproblems-problemphyxmini0407-massdensityquantity, def:physics:phyx-mini-0407:phyxminiproblems-problemphyxmini0407-volumequantity, def:physics:phyx-mini-0407:phyxminiproblems-problemphyxmini0407-accelerationquantity, def:physics:phyx-mini-0407:phyxminiproblems-problemphyxmini0407-temperaturequantity, def:physics:phyx-mini-0407:phyxminiproblems-problemphyxmini0407-molargasconstantquantity, def:physics:phyx-mini-0407:phyxminiproblems-problemphyxmini0407-gasspecies, def:physics:phyx-mini-0407:phyxminiproblems-problemphyxmini0407-pistonmaterial, def:physics:phyx-mini-0407:phyxminiproblems-problemphyxmini0407-cylinderorientation, def:physics:phyx-mini-0407:phyxminiproblems-problemphyxmini0407-abovepistonmedium, def:physics:phyx-mini-0407:phyxminiproblems-problemphyxmini0407-pistonmobility, def:physics:phyx-mini-0407:phyxminiproblems-problemphyxmini0407-verticalmotion, def:physics:phyx-mini-0407:phyxminiproblems-problemphyxmini0407-figureregionlabel}
  The cylinder, piston, gas states, and signed work quantities. workByGas is positive when the gas transfers mechanical energy to the piston. workOnGas is positive when the environment transfers mechanical energy to the gas; the governing laws below relate the two by a minus sign
\end{definition}
\begin{proof}
  This is the declared model definition of Copper Piston Nitrogen Cylinder.
\end{proof}
\begin{definition}[Matches Problem And Primary Figure]
  \label{def:physics:phyx-mini-0407:phyxminiproblems-problemphyxmini0407-matchesproblemandprimaryfigure}
  \lean{PhyXMiniProblems.ProblemPhyXMini0407.MatchesProblemAndPrimaryFigure}
  \uses{def:physics:phyx-mini-0407:phyxminiproblems-problemphyxmini0407-lengthinmeters, def:physics:phyx-mini-0407:phyxminiproblems-problemphyxmini0407-temperatureindegreescelsius, def:physics:phyx-mini-0407:phyxminiproblems-problemphyxmini0407-copperpistonnitrogencylinder}
  The qualitative setup and exact decimal readouts supplied by the prose and primary image. The image's 20 cm arrow runs from the cylinder bottom to the underside of the piston, so it is recorded as the initial gas-column height. No field mentions the final height, rise distance, final temperature, piston density, gravitational acceleration, work, or an answer choice
\end{definition}
\begin{proof}
  This is the declared model definition of Matches Problem And Primary Figure.
\end{proof}
\begin{definition}[Has Physical Parameters]
  \label{def:physics:phyx-mini-0407:phyxminiproblems-problemphyxmini0407-hasphysicalparameters}
  \lean{PhyXMiniProblems.ProblemPhyXMini0407.HasPhysicalParameters}
  \uses{def:physics:phyx-mini-0407:phyxminiproblems-problemphyxmini0407-lengthinmeters, def:physics:phyx-mini-0407:phyxminiproblems-problemphyxmini0407-massinkilograms, def:physics:phyx-mini-0407:phyxminiproblems-problemphyxmini0407-massdensityinkilogramspercubicmeter, def:physics:phyx-mini-0407:phyxminiproblems-problemphyxmini0407-areainsquaremeters, def:physics:phyx-mini-0407:phyxminiproblems-problemphyxmini0407-volumeincubicmeters, def:physics:phyx-mini-0407:phyxminiproblems-problemphyxmini0407-pressureinpascals, def:physics:phyx-mini-0407:phyxminiproblems-problemphyxmini0407-accelerationinmeterspersecondsquared, def:physics:phyx-mini-0407:phyxminiproblems-problemphyxmini0407-temperatureinkelvins, def:physics:phyx-mini-0407:phyxminiproblems-problemphyxmini0407-molargasconstantinsi, def:physics:phyx-mini-0407:phyxminiproblems-problemphyxmini0407-copperpistonnitrogencylinder}
  Positivity assumptions selecting physically meaningful parameters
\end{definition}
\begin{proof}
  This is the declared model definition of Has Physical Parameters.
\end{proof}
\begin{definition}[Obeys Piston Cylinder Laws]
  \label{def:physics:phyx-mini-0407:phyxminiproblems-problemphyxmini0407-obeyspistoncylinderlaws}
  \lean{PhyXMiniProblems.ProblemPhyXMini0407.ObeysPistonCylinderLaws}
  \uses{def:physics:phyx-mini-0407:phyxminiproblems-problemphyxmini0407-lengthinmeters, def:physics:phyx-mini-0407:phyxminiproblems-problemphyxmini0407-massinkilograms, def:physics:phyx-mini-0407:phyxminiproblems-problemphyxmini0407-massdensityinkilogramspercubicmeter, def:physics:phyx-mini-0407:phyxminiproblems-problemphyxmini0407-areainsquaremeters, def:physics:phyx-mini-0407:phyxminiproblems-problemphyxmini0407-volumeincubicmeters, def:physics:phyx-mini-0407:phyxminiproblems-problemphyxmini0407-pressureinpascals, def:physics:phyx-mini-0407:phyxminiproblems-problemphyxmini0407-accelerationinmeterspersecondsquared, def:physics:phyx-mini-0407:phyxminiproblems-problemphyxmini0407-temperatureinkelvins, def:physics:phyx-mini-0407:phyxminiproblems-problemphyxmini0407-molargasconstantinsi, def:physics:phyx-mini-0407:phyxminiproblems-problemphyxmini0407-energyinjoules, def:physics:phyx-mini-0407:phyxminiproblems-problemphyxmini0407-copperpistonnitrogencylinder}
  General geometry, constitutive, equilibrium, ideal-gas, and work laws for the apparatus. The boundary-work field is conditional on constant pressure and therefore states the standard constant-pressure law rather than a numerical answer to this problem. In particular, no field fixes the final height or the requested work
\end{definition}
\begin{proof}
  This is the declared model definition of Obeys Piston Cylinder Laws.
\end{proof}
\begin{definition}[Answer Choice]
  \label{def:physics:phyx-mini-0407:phyxminiproblems-problemphyxmini0407-answerchoice}
  \lean{PhyXMiniProblems.ProblemPhyXMini0407.AnswerChoice}
  Labels of the four answer choices printed with the problem
\end{definition}
\begin{proof}
  This is the declared model definition of Answer Choice.
\end{proof}
\begin{definition}[answer Choice In Joules]
  \label{def:physics:phyx-mini-0407:phyxminiproblems-problemphyxmini0407-answerchoiceinjoules}
  \lean{PhyXMiniProblems.ProblemPhyXMini0407.answerChoiceInJoules}
  \uses{def:physics:phyx-mini-0407:phyxminiproblems-problemphyxmini0407-answerchoice}
  Signed joule values displayed beside the four answer labels
\end{definition}
\begin{proof}
  This is the declared model definition of answer Choice In Joules.
\end{proof}
\begin{definition}[Work On Gas Matches Choice]
  \label{def:physics:phyx-mini-0407:phyxminiproblems-problemphyxmini0407-workongasmatcheschoice}
  \lean{PhyXMiniProblems.ProblemPhyXMini0407.WorkOnGasMatchesChoice}
  \uses{def:physics:phyx-mini-0407:phyxminiproblems-problemphyxmini0407-energyinjoules, def:physics:phyx-mini-0407:phyxminiproblems-problemphyxmini0407-copperpistonnitrogencylinder, def:physics:phyx-mini-0407:phyxminiproblems-problemphyxmini0407-answerchoice, def:physics:phyx-mini-0407:phyxminiproblems-problemphyxmini0407-answerchoiceinjoules}
  A displayed choice matches the signed work-on-gas convention
\end{definition}
\begin{proof}
  This is the declared model definition of Work On Gas Matches Choice.
\end{proof}
\begin{theorem}[problem phyx mini 0407]
  \label{thm:physics:phyx-mini-0407:phyxminiproblems-problemphyxmini0407-problem-phyx-mini-0407}
  \lean{PhyXMiniProblems.ProblemPhyXMini0407.problem_phyx_mini_0407}
  \uses{def:physics:phyx-mini-0407:phyxminiproblems-problemphyxmini0407-lengthinmeters, def:physics:phyx-mini-0407:phyxminiproblems-problemphyxmini0407-massinkilograms, def:physics:phyx-mini-0407:phyxminiproblems-problemphyxmini0407-massdensityinkilogramspercubicmeter, def:physics:phyx-mini-0407:phyxminiproblems-problemphyxmini0407-accelerationinmeterspersecondsquared, def:physics:phyx-mini-0407:phyxminiproblems-problemphyxmini0407-energyinjoules, def:physics:phyx-mini-0407:phyxminiproblems-problemphyxmini0407-copperpistonnitrogencylinder, def:physics:phyx-mini-0407:phyxminiproblems-problemphyxmini0407-matchesproblemandprimaryfigure, def:physics:phyx-mini-0407:phyxminiproblems-problemphyxmini0407-hasphysicalparameters, def:physics:phyx-mini-0407:phyxminiproblems-problemphyxmini0407-obeyspistoncylinderlaws}
  For a rising piston, the work done on the gas is the negative of the piston's gain in gravitational potential energy. The second equality expands the piston mass using its density, circular cross-section, and thickness. The final height remains a genuine variable because the source supplies no rise distance or final temperature. Thus this theorem preserves the physical answer determined by the stated data without falsely selecting a number
\end{theorem}
\begin{proof}
  The conclusion follows from the governing relations and figure readouts named in the statement of problem phyx mini 0407.
\end{proof}
% --- Archon declaration topology end ---
% --- Archon physics formalization source end ---
